% =============================================================================
% chapter4.tex — Глава 4. Тестирование и оценка результатов
% =============================================================================

\section{ТЕСТИРОВАНИЕ И ОЦЕНКА РЕЗУЛЬТАТОВ}

\subsection{Стратегия тестирования}

\todo{Описать пирамиду тестирования:

Общее количество: ~107 тестовых файлов (Kotlin).

Распределение по модулям:
\begin{itemize}
  \item contexts/graph — 40+ тестов (наибольшее покрытие: PSI-парсинг, линковка, экстракторы)
  \item contexts/chunking — 11 тестов (оркестратор, стратегия, NodeDoc, LangGuards)
  \item contexts/rag — 10 тестов (каждый шаг FSM отдельно + интеграция)
  \item kernel/db — 8 тестов (репозитории, интеграция с PostgreSQL)
  \item contexts/git — 7 тестов (GitHub, GitLab, classpath resolver)
  \item contexts/library — 6 тестов (ASM-парсинг, байткод-анализ, интеграционные точки)
  \item contexts/postprocess — 6 тестов (handlers, merger, scheduler)
  \item contexts/ai — 5 тестов (LLM-клиенты: Coder, Talker, Digest, Summary)
  \item contexts/embedding — 3 теста (store, search, config)
\end{itemize}

Фреймворки: JUnit 5 (Jupiter), MockK (Kotlin-native mocking), AssertJ (fluent assertions), Kotlin Test DSL.

Интеграционные тесты: TestContainers (PostgreSQL с pgvector), BaseRepositoryTest (общий контекст для DB-тестов).}

\subsection{Юнит-тестирование ключевых компонентов}

\subsubsection{Тестирование построения графа}

\todo{Описать тесты в contexts/graph (40+ файлов):

KotlinGraphBuilderTest: проверка создания узлов из Kotlin-конструкций:
\begin{itemize}
  \item Парсинг class/interface/enum/data class/object → корректный NodeKind
  \item Извлечение FQN, signature, source code, line span
  \item Обработка generics, suspend fun, extension functions
  \item Фильтрация шума (вспомогательные функции)
\end{itemize}

Тесты линкеров (по одному на каждый линкер):
\begin{itemize}
  \item StructuralEdgeLinkerTest — CONTAINS рёбра
  \item InheritanceEdgeLinkerTest — IMPLEMENTS, INHERITS, EXTENDS
  \item AnnotationEdgeLinkerTest — ANNOTATED\_WITH
  \item CallEdgeLinkerTest — CALLS\_CODE
  \item IntegrationEdgeLinkerTest — CALLS\_HTTP, PRODUCES, CONSUMES
  \item ThrowEdgeLinkerTest — THROWS
\end{itemize}

Тесты экстракторов API-метаданных:
\begin{itemize}
  \item HttpEndpointExtractorTest — парсинг @RequestMapping
  \item GrpcEndpointExtractorTest — gRPC endpoints
  \item MessageBrokerExtractorTest — Kafka/Rabbit
\end{itemize}

Тесты экстракторов NodeKind (15+ файлов): по одному тесту на каждый тип узла.}

\subsubsection{Тестирование чанкования и эмбеддингов}

\todo{Описать тесты:

ChunkBuildOrchestratorImplTest:
\begin{itemize}
  \item Корректное разбиение узлов на чанки
  \item Батчевая обработка (batchSize=50)
  \item Обработка пустого графа
  \item Dry-run режим
\end{itemize}

PerNodeChunkStrategyTest: один узел → один чанк с корректным content и metadata.

ChunkWriterImplTest: персистенция, дедупликация по content\_hash.

EmbeddingSearchServiceImplTest:
\begin{itemize}
  \item searchByText — семантический поиск, проверка similarity score
  \item Пустой запрос → пустые результаты
  \item topK ограничение
\end{itemize}

EmbeddingStoreServiceImplTest: addDocument, deleteDocument, валидация размерности.

NodeDocGeneratorTest: Coder → Talker → Digest цепочка, обработка timeout.

LangGuardsTest: проверка Cyrillic ratio, фильтрация мусорных описаний.}

\subsubsection{Тестирование RAG-системы}

\todo{Описать тесты в contexts/rag (10 файлов):

Каждый шаг FSM тестируется изолированно:
\begin{itemize}
  \item NormalizationStepTest — нормализация запроса
  \item ExtractionStepTest — извлечение сущностей (className, methodName)
  \item ExactSearchStepTest — точный поиск по FQN
  \item GraphExpansionStepTest — обход графа, лимит соседей
  \item RewritingStepTest — переформулировка
  \item ExpansionStepTest — синонимы
  \item VectorSearchStepTest — cosine similarity
  \item RerankingStepTest — сортировка и фильтрация
\end{itemize}

GraphRequestProcessorTest — интеграция FSM: проверка переходов между шагами, обработка timeout, MAX\_ITERATIONS=20.

RagServiceImplTest — полный цикл: ask(query) → context building → LLM call → response formatting.}

\subsubsection{Тестирование кросс-приложенческого анализа}

\todo{Описать 6 тестов CrossAppGraphServiceImplTest:
\begin{enumerate}
  \item Фильтрация по типу интеграции (HTTP only, KAFKA only, CAMEL only)
  \item Группировка по FQN (один endpoint в 2 приложениях → один интеграционный узел)
  \item Определение ролей приложений (producer/consumer для Kafka, client/server для HTTP)
  \item Пустой набор приложений → пустой ответ
  \item Все типы интеграций одновременно
  \item Статистика (applicationCount, httpEndpoints, kafkaTopics, camelRoutes)
\end{enumerate}

Все тесты используют MockK для мокирования NodeRepository.}

\subsection{Интеграционное тестирование}

\todo{Описать интеграционные тесты:

Инфраструктура:
\begin{itemize}
  \item TestContainers: PostgreSQL 16 + pgvector (автоматический запуск Docker-контейнера)
  \item BaseRepositoryTest: общий базовый класс, @Transactional, настройка DataSource
  \item TestApplication: минимальная Spring Boot конфигурация для тестов
\end{itemize}

Репозиторные тесты (kernel/db, 8 файлов):
\begin{itemize}
  \item ApplicationRepositoryTest — CRUD, поиск по key, фильтрация
  \item NodeRepositoryTest — сохранение/поиск узлов, findByFqn, findByKind, findAllByApplicationIdInAndKindIn
  \item EdgeRepositoryTest — составной PK, upsert, каскадное удаление
  \item ChunkRepositoryTest — vector(1024) поля, content\_hash дедупликация
  \item ChunkGraphRepositoryImplTest — графовые запросы через чанки
  \item LibraryRepositoryTest — координаты, фильтрация
  \item NodeDocRepositoryTest — upsert по node\_id, обновление quality
\end{itemize}

Проверка pgvector: INSERT vector(1024), HNSW-индекс, cosine similarity search, корректность nearest-neighbor результатов.}

\subsection{Оценка качества построения графа}

\todo{Описать метрики и провести эксперимент:

Метрики:
\begin{itemize}
  \item Полнота обнаружения узлов: precision/recall по каждому NodeKind (сравнение с ручным подсчётом)
  \item Точность связей: доля корректных рёбер (верификация выборки из 100 рёбер)
  \item Покрытие интеграционных точек: процент обнаруженных HTTP/Kafka/Camel точек vs реальные (по документации/конфигурации)
\end{itemize}

Эксперимент на реальном проекте (Kotlin/Spring Boot):
\begin{itemize}
  \item Входные данные: репозиторий из N .kt файлов, M JAR-зависимостей
  \item Результат: кол-во обнаруженных узлов по типам, кол-во рёбер по типам
  \item Время индексации (clone + parse + link + chunk + embed)
  \item Сравнение с эталонным графом (ручная разметка 50 классов)
\end{itemize}}

\subsection{Оценка качества генерации документации}

\todo{Описать метрики и эксперимент:

Метрики:
\begin{itemize}
  \item Полнота (coverage): процент узлов, для которых сгенерировано описание (docTech + docPublic)
  \item Точность (precision): доля корректных описаний (экспертная оценка выборки из 50 узлов)
  \item Читаемость (readability): оценка понятности docPublic для нетехнического читателя (шкала 1--5, экспертная оценка)
  \item Cyrillic ratio: процент кириллических символов в docPublic (целевое: >=60\%)
  \item Сравнение Coder vs Talker: насколько Talker упрощает техническое описание
\end{itemize}

Сравнение с ручной документацией:
\begin{itemize}
  \item Полнота: процент полей/методов, описанных автоматически vs вручную
  \item Время: автогенерация (секунды) vs ручное написание (минуты/часы)
  \item Актуальность: автодокументация всегда актуальна (генерируется из текущего кода)
\end{itemize}}

\subsection{Оценка RAG-системы}

\todo{Описать метрики и эксперимент:

Метрики:
\begin{itemize}
  \item Релевантность ответов: precision@k, recall@k (набор из 30 вопросов с эталонными ответами)
  \item Качество источников (attribution accuracy): процент ответов, подкреплённых корректными source-ссылками
  \item Полнота контекста: достаточность контекста для ответа (оценка экспертом)
\end{itemize}

Сравнение наивного RAG vs графового расширения:
\begin{itemize}
  \item Наивный RAG: только vector search (без EXACT\_SEARCH и GRAPH\_EXPANSION шагов)
  \item Графовый RAG: полный пайплайн (8 шагов FSM)
  \item Метрика: improvement = (graphRAG\_precision - naiveRAG\_precision) / naiveRAG\_precision
  \item Гипотеза: графовое расширение увеличивает precision на 20--40\% за счёт добавления контекста зависимостей
\end{itemize}

Тестовые вопросы по категориям:
\begin{itemize}
  \item <<Что делает метод X?>> — точное описание (EXTRACTION → EXACT\_SEARCH)
  \item <<Какие сервисы вызывает класс Y?>> — графовые связи (GRAPH\_EXPANSION)
  \item <<Как работает обработка заказов?>> — широкий контекст (VECTOR\_SEARCH + GRAPH\_EXPANSION)
\end{itemize}}

\subsection{Нагрузочное тестирование и отказоустойчивость}

\todo{Описать тестирование:

Circuit Breaker:
\begin{itemize}
  \item Сценарий: отказ Ollama (timeout 30s)
  \item Поведение: после 5 из 10 вызовов с ошибкой → CB открывается → запросы отклоняются 60с → half-open → 5 пробных вызовов → закрытие
  \item Метрика: время восстановления, fallback ответ
\end{itemize}

Rate Limiting:
\begin{itemize}
  \item Сценарий: 150 req/min на /api/v1/rag/ask (лимит: 30 req/min)
  \item Поведение: первые 30 → 200 OK, остальные → 429 Too Many Requests, Retry-After: 60
  \item Проверка: корректность Redis ZSET sliding window
\end{itemize}

Bulkhead:
\begin{itemize}
  \item Сценарий: 10 параллельных LLM-запросов (лимит: 5)
  \item Поведение: первые 5 выполняются, остальные 5 ждут (maxWaitDuration=10s)
  \item Метрика: время ожидания, throughput
\end{itemize}

Графовые операции:
\begin{itemize}
  \item max-nodes=100\,000: проверка OOM-защиты при больших графах
  \item Параллельная линковка: 8 потоков, таймаут 10 мин
\end{itemize}}

\subsection{Сравнение с существующими решениями}

\todo{Сравнительная таблица:

\begin{center}
\begin{tabular}{|l|c|c|c|c|}
\hline
Критерий & Doc-Generator & Copilot Chat & Cody & Наивный RAG \\
\hline
Контекст системы & + & -- & +/-- & -- \\
Интеграции (HTTP/Kafka) & + & -- & -- & -- \\
Конфиденциальность & + (локально) & -- (облако) & +/-- & + \\
Стоимость (inference) & 0 (Ollama) & платно & платно & зависит \\
Графовое расширение & + & -- & -- & -- \\
Кросс-приложенческий & + & -- & -- & -- \\
Offline-работа & + & -- & +/-- & + \\
Качество документации & высокое & среднее & среднее & низкое \\
Покрытие & полное & по запросу & по запросу & частичное \\
\hline
\end{tabular}
\end{center}

Обоснование каждой оценки с примерами. Уникальные преимущества Doc-Generator: графовое расширение контекста, кросс-приложенческий анализ, локальный inference (конфиденциальность), автоматическая актуальность.}

\subsection{Результаты апробации на реальных проектах}

\todo{Описать результаты применения к реальным проектам (Kotlin/Spring Boot):

Проект 1: \textit{(название, размер кодовой базы)}
\begin{itemize}
  \item Входные данные: N .kt файлов, M строк кода, K JAR-зависимостей
  \item Обнаруженные узлы: по типам (CLASS: X, METHOD: Y, ENDPOINT: Z, TOPIC: W, ...)
  \item Обнаруженные рёбра: по типам (CONTAINS: X, CALLS\_CODE: Y, CALLS\_HTTP: Z, ...)
  \item Время индексации: clone (Xs) + parse (Xs) + link (Xs) + chunk (Xs) + embed (Xs) = total
  \item Сгенерировано описаний: docTech (N), docPublic (N), docDigest (N)
  \item Средний Cyrillic ratio: X\%
\end{itemize}

Проект 2: \textit{(аналогично)}

Обратная связь от разработчиков:
\begin{itemize}
  \item Полезность автодокументации (шкала 1--5)
  \item Качество RAG-ответов (шкала 1--5)
  \item Удобство графовой визуализации (шкала 1--5)
  \item Предложения по улучшению
\end{itemize}}
